\documentclass[12pt]{article}
\usepackage[brazil]{babel}
\usepackage[utf8]{inputenc}
\usepackage[T1]{fontenc}
\usepackage{amsmath, amssymb}
\usepackage{geometry}
\usepackage{graphicx}
\usepackage{tikz}
\usetikzlibrary{shapes,arrows,positioning}
\usepackage{xcolor}
\usepackage{enumitem}
\usepackage{caption}
\usepackage{float}

\definecolor{maincolor}{RGB}{0,100,50}
\usepackage{fancyhdr}
\definecolor{forestgreen}{RGB}{0,100,50}
\usepackage[most]{tcolorbox}

\geometry{margin=2.5cm}
\pagestyle{fancy}
\fancyhf{}

\renewcommand{\headrulewidth}{2pt}
\renewcommand{\footrulewidth}{2pt}

\renewcommand{\headrule}{\hbox to\headwidth{%
  \color{forestgreen}\leaders\hrule height \headrulewidth\hfill}}

\renewcommand{\footrule}{\hbox to\headwidth{%
  \color{forestgreen}\leaders\hrule height \footrulewidth\hfill}}

\setlength{\headheight}{51pt}

\fancyhead[L]{%
  \includegraphics[height=1.2cm]{Universidade_de_Rio_Verde.jpg}
}

\fancyhead[R]{%
  \raggedleft
  \textcolor{forestgreen}{%
    \small
    \textbf{Lista de Exercícios — Funções}\\
    Me.~Sergio Souza Novak\\
    \textit{Lógica para Engenharia de Software}
  } 
}

\fancyfoot[C]{%
  \textcolor{forestgreen}{\thepage}
}

\begin{document}

\begin{enumerate}[label=\arabic*), leftmargin=1.5cm, itemsep=1cm]

    \item Determine o domínio de definição das seguintes funções reais, indicando as restrições necessárias em cada caso:
    \begin{enumerate}[label=\alph*)]
        \item $f(x) = \dfrac{x + 5}{x - 7}$
        \item $g(x) = \sqrt{2x - 10}$
        \item $h(x) = \ln(x^2 - 16)$
        \item $i(x) = \dfrac{\sqrt{x + 2}}{x - 3}$
        \item $j(x) = \sqrt[3]{x - 8}$
        \item $k(x) = \dfrac{1}{x^2 - 5x + 6}$
        \item $l(x) = \log(x - 2) + \sqrt{10 - x}$
        \item $m(x) = \dfrac{1}{\sqrt{x^2 - 9}}$
        \item $n(x) = \sqrt{\dfrac{x-1}{x+2}}$
    \end{enumerate}

    \item Dadas as funções $f(x) = 2x - 3$ e $g(x) = x^2 + 4$, determine as expressões para as funções compostas abaixo e simplifique o resultado:
    \begin{enumerate}[label=\alph*)]
        \item $(f \circ g)(x)$
        \item $(g \circ f)(x)$
        \item $(f \circ f)(x)$
        \item Calcule o valor numérico de $(f \circ g)(2)$.
    \end{enumerate}

    \item Analise as funções abaixo e classifique-as em Injetora, Sobrejetora ou Bijetora. Justifique sua resposta analisando o domínio ($D$) e o contradomínio ($CD$).
    \begin{enumerate}[label=\alph*)]
        \item $f: \mathbb{R} \to \mathbb{R}, \quad f(x) = 3x + 5$
        \item $g: \mathbb{R} \to \mathbb{R}, \quad g(x) = x^2$
        \item $h: \mathbb{R} \to [0, \infty), \quad h(x) = x^2$
        \item $i: [0, \infty) \to [0, \infty), \quad i(x) = x^2$
        \item $j: \mathbb{R} \to \mathbb{R}, \quad j(x) = x^3$
        \item $k: \mathbb{R} \to \mathbb{R}, \quad k(x) = |x|$
        \item $l: \mathbb{R} \to (0, \infty), \quad l(x) = e^x$
        \item $m: (0, \infty) \to \mathbb{R}, \quad m(x) = \ln(x)$
        \item $n: \mathbb{R} \setminus \{0\} \to \mathbb{R} \setminus \{0\}, \quad n(x) = \frac{1}{x}$
        \item $p: \mathbb{R} \to \mathbb{R}, \quad p(x) = \sin(x)$
    \end{enumerate}


    \item Utilize o Teste da Reta Vertical e o Teste da Reta Horizontal para classificar as curvas representadas nos gráficos abaixo como: (I) Função ou não-função; (II) Injetora ou não-injetora.

    \begin{center}
    \begin{tikzpicture}[scale=0.6]
        % Linha 1
        % Gráfico A: Parábola (Função, Não Injetora)
        \begin{scope}[xshift=0cm, yshift=0cm]
            \draw[->] (-2,0) -- (2,0) node[right] {$x$};
            \draw[->] (0,-1) -- (0,3) node[above] {$y$};
            \draw[domain=-1.5:1.5, smooth, blue, thick] plot ({\x}, {\x*\x});
            \node at (0,-1.5) {Gráfico A};
        \end{scope}

        % Gráfico B: Parábola Lateral (Não Função)
        \begin{scope}[xshift=6cm, yshift=0cm]
            \draw[->] (-1,0) -- (3,0) node[right] {$x$};
            \draw[->] (0,-2) -- (0,2) node[above] {$y$};
            \draw[domain=-1.5:1.5, smooth, blue, thick] plot ({\x*\x}, {\x});
            \node at (1,-2.5) {Gráfico B};
        \end{scope}

        % Gráfico C: Reta Inclinada (Função, Injetora)
        \begin{scope}[xshift=12cm, yshift=0cm]
            \draw[->] (-2,0) -- (2,0) node[right] {$x$};
            \draw[->] (0,-1) -- (0,3) node[above] {$y$};
            \draw[domain=-1.5:1.5, smooth, blue, thick] plot ({\x}, {\x + 1});
            \node at (0,-1.5) {Gráfico C};
        \end{scope}

        % Linha 2
        % Gráfico D: Cúbica (Função, Injetora)
        \begin{scope}[xshift=0cm, yshift=-6cm]
            \draw[->] (-2,0) -- (2,0) node[right] {$x$};
            \draw[->] (0,-2) -- (0,2) node[above] {$y$};
            \draw[domain=-1.2:1.2, smooth, blue, thick] plot ({\x}, {\x*\x*\x});
            \node at (0,-2.5) {Gráfico D};
        \end{scope}

        % Gráfico E: Círculo (Não Função)
        \begin{scope}[xshift=6cm, yshift=-6cm]
            \draw[->] (-2,0) -- (2,0) node[right] {$x$};
            \draw[->] (0,-2) -- (0,2) node[above] {$y$};
            \draw[blue, thick] (0,0) circle (1.2);
            \node at (0,-2.5) {Gráfico E};
        \end{scope}

        % Gráfico F: Reta Horizontal (Função, Não Injetora)
        \begin{scope}[xshift=12cm, yshift=-6cm]
            \draw[->] (-2,0) -- (2,0) node[right] {$x$};
            \draw[->] (0,-1) -- (0,3) node[above] {$y$};
            \draw[blue, thick] (-1.8, 1.5) -- (1.8, 1.5);
            \node at (0,-1.5) {Gráfico F};
        \end{scope}

        % Linha 3
        % Gráfico G: Senoide (Função, Não Injetora)
        \begin{scope}[xshift=0cm, yshift=-12cm]
            \draw[->] (-2,0) -- (2,0) node[right] {$x$};
            \draw[->] (0,-2) -- (0,2) node[above] {$y$};
            \draw[domain=-1.8:1.8, smooth, blue, thick] plot ({\x}, {sin(\x*180)});
            \node at (0,-2.5) {Gráfico G};
        \end{scope}

        % Gráfico H: V (Valor Absoluto) (Função, Não Injetora)
        \begin{scope}[xshift=6cm, yshift=-12cm]
            \draw[->] (-2,0) -- (2,0) node[right] {$x$};
            \draw[->] (0,-1) -- (0,3) node[above] {$y$};
            \draw[blue, thick] (-1.5, 1.5) -- (0,0) -- (1.5, 1.5);
            \node at (0,-1.5) {Gráfico H};
        \end{scope}

        % Gráfico I: Semicírculo Superior (Função, Não Injetora)
        \begin{scope}[xshift=12cm, yshift=-12cm]
            \draw[->] (-2,0) -- (2,0) node[right] {$x$};
            \draw[->] (0,-1) -- (0,3) node[above] {$y$};
            \draw[domain=0:180, smooth, blue, thick] plot ({1.5*cos(\x)}, {1.5*sin(\x)});
            \node at (0,-1.5) {Gráfico I};
        \end{scope}

        % Linha 4
        % Gráfico J: Exponencial (Função, Injetora)
        \begin{scope}[xshift=3cm, yshift=-18cm]
            \draw[->] (-2,0) -- (2,0) node[right] {$x$};
            \draw[->] (0,-1) -- (0,3) node[above] {$y$};
            \draw[domain=-1.8:1.1, smooth, blue, thick] plot ({\x}, {exp(\x)});
            \node at (0,-1.5) {Gráfico J};
        \end{scope}

        % Gráfico K: Logaritmo (Função, Injetora)
        \begin{scope}[xshift=9cm, yshift=-18cm]
            \draw[->] (-1,0) -- (3,0) node[right] {$x$};
            \draw[->] (0,-2) -- (0,2) node[above] {$y$};
            \draw[domain=0.15:2.5, smooth, blue, thick] plot ({\x}, {ln(\x)});
            \node at (1,-2.5) {Gráfico K};
        \end{scope}
    \end{tikzpicture}
    \end{center}

    \item Considere a função $f: A \to B$ representada pelo diagrama de Venn abaixo, onde $A = \{1, 2, 3\}$ e $B = \{a, b, c, d\}$.
    \begin{center}
    \begin{tikzpicture}[scale=1]
        \node (1) at (0,1) {1};
        \node (2) at (0,0) {2};
        \node (3) at (0,-1) {3};
        \draw (0,0) ellipse (0.8 and 1.8);
        \node at (0,2) {$A$};

        \node (a) at (3,1.5) {a};
        \node (b) at (3,0.5) {b};
        \node (c) at (3,-0.5) {c};
        \node (d) at (3,-1.5) {d};
        \draw (3,0) ellipse (0.8 and 2);
        \node at (3,2.3) {$B$};

        \draw[->, >=stealth] (1) -- (a);
        \draw[->, >=stealth] (2) -- (c);
        \draw[->, >=stealth] (3) -- (c);
    \end{tikzpicture}
    \end{center}

    Determine:
    \begin{enumerate}[label=\alph*)]
        \item O conjunto Domínio ($D$).
        \item O conjunto Contradomínio ($CD$).
        \item O conjunto Imagem ($Im$).
        \item Esta função é sobrejetora? Justifique.
        \item Esta função é injetora? Justifique.
    \end{enumerate}

\end{enumerate}

\end{document}
